\section{讨论与结论}

实验结果表明，在区域风电集群 24 h 预测中，未来气象预报仍是主导信息来源。Weather-only 已经强于多种网格气象编码基线，说明 Pangu-Weather 网格化预报能够捕捉预测窗口内的主要外生驱动。本文方法的作用不是用历史功率替代未来气象，而是在气象驱动预测中引入预测起点前的真实场群状态约束，使模型能够在读取未来气象时同时感知当前运行状态。

逐预测步结果进一步说明，运行状态约束的收益具有明确的业务含义：它在完整 24 h 平均指标上带来稳定改善，并在目标窗口起始时段表现最为突出。H1 MAE 的大幅降低表明，当前出力水平、短期惯性和多场站同步变化能够有效修正气象到功率映射的初始偏移；同时，24 步平均 MAE 和 RMSE 的同步下降说明这种约束会通过 query 解码过程影响完整预测序列，而不是只作用于单个时刻。

消融实验进一步区分了两类信息来源。No wind 和 No memory 的明显退化表明，风场变量和时空气象记忆是模型的主信息载体；Random power 和 Zero power 的退化则表明，功率侧收益来自真实运行状态，而不是额外分支或参数量。由此可见，本文方法的关键在于非对称融合：未来气象提供目标窗口内的主驱动，近期场站功率提供预测起点的运行状态约束，二者在 memory-level fusion 中共同服务于区域风电预测。

本文提出了运行状态约束的气象-功率融合网络，将未来气象编码为主记忆，将近期场站级功率编码为动态上下文 token，并通过 memory-level fusion 条件化 future queries 的读取过程。吉林区域 58 个风电场站案例验证了该方法在整体预测误差和起始时段预测误差上的有效性。未来工作可进一步结合气象预报不确定性、跨区域迁移和极端爬坡事件建模，提升模型在复杂运行场景下的适应性。
